% ============================================================
%  Assignment 1 — Satisfiability & SAT Solvers
%  Beautiful LaTeX Report
% ============================================================
\documentclass[11pt, a4paper]{article}

% ── Core packages ────────────────────────────────────────────
\usepackage[a4paper, top=2.5cm, bottom=2.5cm, left=2.8cm, right=2.8cm, headheight=15pt]{geometry}
\usepackage[T1]{fontenc}
\usepackage[utf8]{inputenc}

\usepackage{verbatim}
% ── Color scheme (matching the blue theme) ───────────────────
\usepackage[dvipsnames, table]{xcolor}
\definecolor{primaryblue}{RGB}{31, 73, 125}      % deep navy
\definecolor{accentblue}{RGB}{46, 116, 181}      % medium blue
\definecolor{lightblue}{RGB}{214, 234, 248}      % very light blue
\definecolor{infobox}{RGB}{235, 243, 251}        % info box bg
\definecolor{codebg}{RGB}{245, 245, 245}         % code background
\definecolor{ruleblue}{RGB}{46, 116, 181}        % rule / line color
\definecolor{darkgray}{RGB}{64, 64, 64}
\definecolor{captioncolor}{RGB}{96, 96, 96}
\definecolor{greenbox}{RGB}{226, 239, 218}
\definecolor{greenborder}{RGB}{84, 130, 53}

% ── Typography ────────────────────────────────────────────────
\usepackage{titlesec}
\usepackage{parskip}

% Section styling
\titleformat{\section}
  {\Large\bfseries\color{primaryblue}}
  {\thesection.}{0.6em}{}
  [\vspace{-0.4em}{\color{ruleblue}\rule{\linewidth}{1.2pt}}\vspace{0.3em}]

\titleformat{\subsection}
  {\normalsize\bfseries\color{accentblue}}
  {\thesubsection}{0.6em}{}

\titleformat{\subsubsection}
  {\normalsize\bfseries\color{darkgray}}
  {\thesubsubsection}{0.6em}{}

\titlespacing*{\section}{0pt}{1.4em}{0.5em}
\titlespacing*{\subsection}{0pt}{1em}{0.3em}

% ── Headers & Footers ────────────────────────────────────────
\usepackage{fancyhdr}
\pagestyle{fancy}
\fancyhf{}
% \fancyhead[L]{\small\color{accentblue}\textbf{Assignment 1} — Satisfiability \& SAT Solvers}
% \fancyhead[R]{\small\color{accentblue}\thepage}
% \fancyfoot[C]{\small\color{captioncolor}Submitted: 10 March 2026 \quad|\quad Seminar: 11 March 2026}
\renewcommand{\headrulewidth}{0.5pt}
\renewcommand{\headrule}{\hbox to\headwidth{\color{ruleblue}\leaders\hrule height \headrulewidth\hfill}}
\renewcommand{\footrulewidth}{0pt}

% ── Math & Logic ─────────────────────────────────────────────
\usepackage{amsmath, amssymb, amsthm}
\usepackage{mathtools}

% ── Algorithms ───────────────────────────────────────────────
\usepackage{algorithm}
\usepackage{algpseudocode}
\algnewcommand\algorithmicforeach{\textbf{for each}}
\algdef{S}[FOR]{ForEach}[1]{\algorithmicforeach\ #1\ \algorithmicdo}
\renewcommand{\algorithmicrequire}{\textbf{Input:}}
\renewcommand{\algorithmicensure}{\textbf{Output:}}

% ── Code listings ────────────────────────────────────────────
\usepackage{listings}
\lstdefinestyle{pythonstyle}{
  backgroundcolor=\color{codebg},
  basicstyle=\ttfamily\footnotesize\color{black},
  keywordstyle=\bfseries\color{primaryblue},
  commentstyle=\itshape\color{captioncolor},
  stringstyle=\color{accentblue},
  numberstyle=\tiny\color{captioncolor},
  numbers=left,
  numbersep=8pt,
  stepnumber=1,
  breaklines=true,
  frame=single,
  framerule=0pt,
  framexleftmargin=10pt,
  xleftmargin=20pt,
  showstringspaces=false,
  tabsize=4,
  language=Python,
  morekeywords={Bool, And, Not, Solver, sat, True, False},
  rulecolor=\color{ruleblue},
  belowskip=0.5em,
  aboveskip=0.5em,
}
\lstdefinestyle{inlinestyle}{
  backgroundcolor=\color{codebg},
  basicstyle=\ttfamily\small,
  frame=none,
  breaklines=true,
  showstringspaces=false,
}

% ── Colored boxes ────────────────────────────────────────────
\usepackage[most]{tcolorbox}

\tcbset{
  infostyle/.style={
    colback=infobox,
    colframe=accentblue,
    boxrule=1pt,
    arc=4pt,
    left=8pt, right=8pt, top=6pt, bottom=6pt,
    fontupper=\small\color{primaryblue},
  },
  keystyle/.style={
    colback=lightblue,
    colframe=primaryblue,
    boxrule=1.2pt,
    arc=4pt,
    left=8pt, right=8pt, top=6pt, bottom=6pt,
    fontupper=\small,
  },
  greenstyle/.style={
    colback=greenbox,
    colframe=greenborder,
    boxrule=1pt,
    arc=4pt,
    left=8pt, right=8pt, top=6pt, bottom=6pt,
    fontupper=\small,
  },
  codebox/.style={
    colback=codebg,
    colframe=ruleblue!40,
    boxrule=0.6pt,
    arc=3pt,
    left=6pt, right=6pt, top=4pt, bottom=4pt,
    fontupper=\ttfamily\small,
  },
}

% ── Tables ───────────────────────────────────────────────────
\usepackage{booktabs}
\usepackage{array}
\usepackage{tabularx}

% ── Graphics & Figures ───────────────────────────────────────
\usepackage{graphicx}
\usepackage{caption}
\captionsetup{
  font={small, it},
  labelfont={bf, color=accentblue},
  textfont={color=captioncolor},
  labelsep=period,
  skip=4pt,
}

% ── Hyperlinks ───────────────────────────────────────────────
\usepackage[colorlinks=true, linkcolor=accentblue, urlcolor=accentblue, citecolor=accentblue]{hyperref}

% ── Enumitem ─────────────────────────────────────────────────
\usepackage{enumitem}
\setlist[itemize]{leftmargin=1.4em, itemsep=2pt, topsep=4pt}
\setlist[enumerate]{leftmargin=1.4em, itemsep=2pt, topsep=4pt}

% ══════════════════════════════════════════════════════════════
\begin{document}
% ══════════════════════════════════════════════════════════════

% ── Title Page ───────────────────────────────────────────────
\begin{titlepage}
\centering
\vspace*{3cm}

{\color{primaryblue}\rule{\linewidth}{2.5pt}}
\vspace{0.6em}

{\Huge\bfseries\color{primaryblue} Assignment 1}\\[0.4em]
{\Large\color{accentblue} Satisfiability \& SAT Solvers}

\vspace{0.5em}
{\color{primaryblue}\rule{\linewidth}{2.5pt}}

\vspace{1.5cm}

\begin{tcolorbox}[keystyle, width=0.82\linewidth]
\centering
\textbf{Contents of this Report}\\[0.5em]
\begin{itemize}
\item \textbf{Exercise 1:} The DPLL Algorithm — CNF, Inference Rules, Pseudocode \& Worked Example\\[0.4em]
\item \textbf{Exercise 2:} Z3 SMT Solver — Installation, Formula Analysis \& Python Implementation
\end{itemize}
\end{tcolorbox}

\vspace{2cm}

% \begin{tabular}{rl}
% \textbf{\color{accentblue}Submission Deadline:} & 10 March 2026 \\[0.3em]
% \textbf{\color{accentblue}Seminar Date:}         & 11 March 2026 \\[0.3em]
% \textbf{\color{accentblue}Group Size:}           & Up to 3 students \\[0.3em]
% \end{tabular}

\vfill

{\small\color{captioncolor}\itshape Formal Methods — Logic \& Verification}

\end{titlepage}

% ── Table of Contents ────────────────────────────────────────
\tableofcontents
\newpage

% ══════════════════════════════════════════════════════════════
\section{Propositional Logic \& Conjunctive Normal Form}
% ══════════════════════════════════════════════════════════════

Before understanding the DPLL algorithm, it is essential to grasp
\emph{Conjunctive Normal Form} (CNF), the canonical representation on which
DPLL operates.

% ── 1.1 ──────────────────────────────────────────────────────
\subsection{Propositional Formulas}

A propositional formula is built from Boolean variables (e.g.\ $p, q, r$)
and the logical connectives $\lnot$, $\land$, $\lor$, $\rightarrow$,
$\leftrightarrow$.
A \textbf{satisfying assignment} (or \emph{model}) is a mapping of variables
to $\{\textit{True},\textit{False}\}$ that makes the formula evaluate to
$\textit{True}$.

% ── 1.2 ──────────────────────────────────────────────────────
\subsection{Conjunctive Normal Form (CNF)}

\begin{tcolorbox}[keystyle]
A formula is in \textbf{CNF} if it is a \emph{conjunction} (\textbf{AND}) of
\emph{clauses}, where every clause is a \emph{disjunction} (\textbf{OR}) of
\emph{literals}.
A \textbf{literal} is a variable $x$ or its negation $\lnot x$.
\end{tcolorbox}

\vspace{0.5em}
Formally:
\[
  \underbrace{(l_{1,1} \lor \cdots \lor l_{1,k_1})}_{\text{clause }1}
  \;\land\;
  \underbrace{(l_{2,1} \lor \cdots \lor l_{2,k_2})}_{\text{clause }2}
  \;\land\;\cdots\;\land\;
  \underbrace{(l_{n,1} \lor \cdots \lor l_{n,k_n})}_{\text{clause }n}
\]

\textbf{Example CNF formula:}
\[
  (A \lor B)\;\land\;(\lnot A \lor C)\;\land\;(\lnot B \lor \lnot C)
\]

% ── 1.3 ──────────────────────────────────────────────────────
\subsection{Converting Any Formula to CNF}

Any propositional formula can be mechanically converted to CNF:

\begin{itemize}
  \item \textbf{Step 1 — Eliminate implications:}
        $A \to B \;\equiv\; \lnot A \lor B$
  \item \textbf{Step 2 — Push negations inward} (De Morgan's Laws):
        $\lnot(A \land B) \equiv \lnot A \lor \lnot B$, \quad
        $\lnot(A \lor B) \equiv \lnot A \land \lnot B$
  \item \textbf{Step 3 — Eliminate double negations:}
        $\lnot\lnot A \equiv A$
  \item \textbf{Step 4 — Distribute $\lor$ over $\land$:}
        $A \lor (B \land C) \equiv (A \lor B) \land (A \lor C)$
\end{itemize}

After these steps the formula is in CNF and ready for DPLL.

% ══════════════════════════════════════════════════════════════
\section{The DPLL Algorithm}
% ══════════════════════════════════════════════════════════════

The \textbf{Davis–Putnam–Logemann–Loveland (DPLL)} algorithm (1962) is a
complete, backtracking-based procedure for deciding the satisfiability of
CNF formulas.  It underpins virtually all modern SAT solvers.

% ── 2.1 ──────────────────────────────────────────────────────
\subsection{Core Inference Rules}

DPLL applies three simplification strategies before any branching.

% ─────────────────────────────────────────────────────────────
\subsubsection{Unit Propagation}

\begin{tcolorbox}[infostyle]
\textbf{Rule:} If a clause contains exactly one literal $\ell$
(a \emph{unit clause}), then $\ell$ must be assigned $\textit{True}$.
All clauses containing $\ell$ are removed; $\lnot\ell$ is removed from every
remaining clause.  Repeat until no unit clause exists.
\end{tcolorbox}

\textbf{Example:}
\[
  (p)\;\land\;(\lnot p \lor q)\;\land\;(\lnot q \lor r)
  \xrightarrow{p=\top}
  (q)\;\land\;(\lnot q \lor r)
  \xrightarrow{q=\top}
  (r)
  \xrightarrow{r=\top}
  \emptyset \;\Rightarrow\; \textbf{SAT}
\]

% ─────────────────────────────────────────────────────────────
\subsubsection{Pure Literal Elimination}

\begin{tcolorbox}[infostyle]
\textbf{Rule:} A literal $\ell$ is \emph{pure} if $\lnot\ell$ never appears
in the formula.  Assigning $\ell = \textit{True}$ satisfies every clause
containing it without breaking any other clause.
\end{tcolorbox}

\textbf{Example:}
$(A \lor B)\;\land\;(\lnot B \lor C)\;\land\;(A \lor \lnot C)$
\;—\; $A$ is pure $\Rightarrow$ assign $A = \top$, remove clauses
containing $A$; result: $(\lnot B \lor C)$.

% ─────────────────────────────────────────────────────────────
\subsubsection{Backtracking (Case Split)}

\begin{tcolorbox}[infostyle]
\textbf{Rule:} When neither UP nor PLE applies, choose an unassigned variable
$x$ and recurse on two branches: $x = \top$ and $x = \bot$.
If both branches fail, the formula is \textbf{UNSAT}.
\end{tcolorbox}

% ── 2.2 ──────────────────────────────────────────────────────
\subsection{DPLL Pseudocode}
\begin{algorithm}[H]
\caption{DPLL($F$, $\alpha$) — Decide satisfiability of CNF formula $F$}
\begin{algorithmic}[1]
\Require CNF formula $F$, partial assignment $\alpha$
\Ensure \textbf{SAT} with model, or \textbf{UNSAT}
\Statex
\Comment{--- Unit Propagation ---}
\While{$\exists$ unit clause $(\ell)$ in $F$}
  \State $F \gets \mathrm{simplify}(F, \ell)$
  \State $\alpha \gets \alpha \cup \{\mathrm{var}(\ell) \mapsto \mathrm{pol}(\ell)\}$
  \If{$\square \in F$}
    \Return \textbf{UNSAT}
  \EndIf
\EndWhile
\Statex
\Comment{--- Pure Literal Elimination ---}
\ForEach{pure literal $\ell$ in $F$}
  \State $F \gets \mathrm{simplify}(F, \ell)$
  \State $\alpha \gets \alpha \cup \{\mathrm{var}(\ell) \mapsto \mathrm{pol}(\ell)\}$
\EndFor
\Statex
\Comment{--- Base Cases ---}
\If{$F = \emptyset$}
  \Return \textbf{SAT}, $\alpha$
\EndIf
\If{$\square \in F$}
  \Return \textbf{UNSAT}
\EndIf
\Statex
\Comment{--- Branching ---}
\State $x \gets \mathrm{chooseVariable}(F)$
\If{$\mathrm{DPLL}(\mathrm{simplify}(F,\, x),\; \alpha \cup \{x \mapsto \top\}) = \textbf{SAT}$}
  \Return \textbf{SAT}
\EndIf
\If{$\mathrm{DPLL}(\mathrm{simplify}(F,\,\lnot x),\; \alpha \cup \{x \mapsto \bot\}) = \textbf{SAT}$}
  \Return \textbf{SAT}
\EndIf
\Return \textbf{UNSAT}
\end{algorithmic}
\end{algorithm}
% ── 2.3 ──────────────────────────────────────────────────────
\subsection{Worked Example}

We run DPLL on $\varphi = (A \lor B)\;\land\;(\lnot A \lor C)\;\land\;(\lnot B \lor \lnot C)$.

\medskip
\noindent\textbf{Initial state.}
No unit clauses; no pure literals.
Choose variable $A$.

\medskip
\noindent\textbf{Branch $A = \top$:}

\[
  \{A,B\},\;\{\lnot A, C\},\;\{\lnot B,\lnot C\}
  \xrightarrow{A=\top}
  \{C\},\;\{\lnot B,\lnot C\}
  \xrightarrow{\text{UP: }C=\top}
  \{\lnot B\}
  \xrightarrow{\text{UP: }B=\bot}
  \emptyset
\]

\begin{tcolorbox}[greenstyle]
\textbf{SAT} \quad Model: $\;A = \top,\; C = \top,\; B = \bot$
\end{tcolorbox}

\noindent The second branch is never reached because the first already yields
a model.

% ── 2.4 ──────────────────────────────────────────────────────
\subsection{Illustrations}

\begin{figure}[H]
\centering
\includegraphics[width=\linewidth]{fig1.png}
\caption{DPLL search-tree execution source - Wikipedia.
Each node represents a partial assignment; leaves marked \textbf{SAT}
or \textbf{UNSAT} are terminal states reached after Unit Propagation.}
\end{figure}

% \begin{figure}[H]
% \centering
% \includegraphics[width=0.70\linewidth]{fig2.png}
% \caption{DAG representation used by the linear solver approach.
% Syntactically identical subformulas are shared as single nodes,
% reducing redundancy.  The marking algorithm checks consistency locally
% at each node, achieving $O(|V|+|E|)$ time complexity.}
% \end{figure}

% ── 2.5 ──────────────────────────────────────────────────────
\subsection{Completeness \& Extension to CDCL}

DPLL is \emph{complete}: it always terminates and correctly reports SAT or
UNSAT.  Modern solvers extend DPLL with
\textbf{Conflict-Driven Clause Learning (CDCL)}: when a contradiction is
found, the conflict graph is analysed and a \emph{learned clause} is added to
prune future search, yielding dramatic practical speed-ups.  Z3 uses a CDCL
engine at its core.

% ══════════════════════════════════════════════════════════════
\section{Exercise 2 — Z3 SMT Solver}
% ══════════════════════════════════════════════════════════════

% ── 3.1 ──────────────────────────────────────────────────────
\subsection{What is Z3?}

Z3 is an open-source \textbf{Satisfiability Modulo Theories (SMT)} solver
developed by Microsoft Research.
It extends CDCL-based SAT solving with support for theories such as linear
arithmetic, arrays, bit-vectors, and uninterpreted functions.
For pure propositional logic, Z3 uses its CDCL engine directly.

% ── 3.2 ──────────────────────────────────────────────────────
\subsection{Installation}

\begin{tcolorbox}[codebox]
\$ pip install z3-solver 
\end{tcolorbox}

% ── 3.3 ──────────────────────────────────────────────────────
\subsection{The Formula: $p \land \lnot\lnot(\lnot q \land \lnot\lnot p)$}

We simplify the formula manually before coding it.

\begin{align*}
  &p \land \lnot\lnot(\lnot q \land \lnot\lnot p) \\
  \equiv\;& p \land \lnot\lnot(\lnot q \land p)
            &&\text{(eliminate }\lnot\lnot p\text{)}\\
  \equiv\;& p \land (\lnot q \land p)
            &&\text{(eliminate outer }\lnot\lnot\text{)}\\
  \equiv\;& p \land \lnot q \land p
            &&\text{(expand)}\\
  \equiv\;& p \land \lnot q
            &&\text{(idempotency: }p\land p \equiv p\text{)}
\end{align*}

\textbf{CNF:} $(p)\;\land\;(\lnot q)$

\begin{tcolorbox}[keystyle]
\textbf{Unique satisfying model:} $p = \textit{True},\quad q = \textit{False}$\\[0.3em]
\textbf{Verification:}\;
$\top \;\land\; \lnot\bot = \top \;\land\; \top = \top$ \checkmark
\end{tcolorbox}

% ── 3.4 ──────────────────────────────────────────────────────
\subsection{Python Implementation}

\begin{lstlisting}[style=pythonstyle]
from z3 import Bool, And, Not, Solver, sat

# Declare propositional variables
p = Bool('p')
q = Bool('q')

# Build formula: p /\ ~~(~q /\ ~~p)
formula = And(
    p,
    Not(Not(And(
        Not(q),
        Not(Not(p))
    )))
)

# Create solver and check
s = Solver()
s.add(formula)
result = s.check()

if result == sat:
    model = s.model()
    print("Result : SAT")
    print("Model  :", model)
    for decl in model.decls():
        print(f"  {decl.name()} = {model[decl]}")
else:
    print("Result : UNSAT")
\end{lstlisting}

% ── 3.5 ──────────────────────────────────────────────────────
\subsection{Expected Output}

\begin{tcolorbox}[codebox]
Result : SAT\\
Model  : [q = False, p = True]\\
\phantom{Model  : }q = False\\
\phantom{Model  : }p = True
\end{tcolorbox}

% ── 3.6 ──────────────────────────────────────────────────────
\subsection{Manual Verification}

Given $p = \top,\; q = \bot$:

\medskip
\begin{tabular}{@{}lll@{}}
  \toprule
  \textbf{Step} & \textbf{Sub-expression} & \textbf{Value} \\
  \midrule
  1 & $\lnot q$             & $\lnot\bot = \top$ \\
  2 & $\lnot\lnot p$        & $\lnot\lnot\top = \top$ \\
  3 & $\lnot q \land \lnot\lnot p$ & $\top \land \top = \top$ \\
  4 & $\lnot\lnot(\cdots)$  & $\top$ \\
  5 & $p \land \top$        & $\top \land \top = \top$ \checkmark \\
  \bottomrule
\end{tabular}

% ── 3.7 ──────────────────────────────────────────────────────
\subsection{Algorithm for DPLL-based Z3 Resolution}

The following algorithm captures how Z3 internally resolves the given formula
using its CDCL (DPLL-based) engine before applying theory reasoning.

\begin{algorithm}[H]
\caption{Z3 CDCL Resolution of $p \land \lnot\lnot(\lnot q \land \lnot\lnot p)$}
\begin{algorithmic}[1]
\Require Formula $\varphi = p \land \lnot\lnot(\lnot q \land \lnot\lnot p)$
\Ensure SAT with model, or UNSAT
\State $\varphi' \gets \mathrm{toCNF}(\varphi)$
       \Comment{Apply double-neg elim \& simplification $\Rightarrow$ $(p)\land(\lnot q)$}
\State $\alpha \gets \{\}$
       \Comment{Empty assignment}
\While{$\exists$ unit clause $(\ell)$ in $\varphi'$}
  \State $\alpha \gets \alpha \cup \{\mathrm{var}(\ell) \mapsto \mathrm{pol}(\ell)\}$
  \State $\varphi' \gets \mathrm{simplify}(\varphi', \ell)$
\EndWhile
\Comment{UP forces $p = \top$ from clause $(p)$, then $q = \bot$ from $(\lnot q)$}
\If{$\varphi' = \emptyset$}
  \Return \textbf{SAT}, $\alpha = \{p \mapsto \top,\; q \mapsto \bot\}$
\EndIf
\Return \textbf{UNSAT}
\end{algorithmic}
\end{algorithm}

% ── 3.8 ──────────────────────────────────────────────────────
\subsection{Comparison: Linear Solver vs.\ DPLL vs.\ Z3}

\begin{table}[H]
\centering
\setlength{\tabcolsep}{8pt}
\renewcommand{\arraystretch}{1.3}
\begin{tabularx}{0.92\linewidth}{>{\bfseries}l X X X}
\toprule
\rowcolor{primaryblue}
{\color{white}Approach} &
{\color{white}Supports $\lor$?} &
{\color{white}Complexity} &
{\color{white}Backtracking?} \\
\midrule
Linear Solver & No (Horn / DAG only) & $O(|V|+|E|)$ & No \\
\rowcolor{lightblue}
DPLL & Yes (general CNF) & Exponential (worst) & Yes \\
Z3 (CDCL+SMT) & Yes + theories & Exp.\ (CDCL pruning) & Yes + learning \\
\bottomrule
\end{tabularx}
\caption{Summary of satisfiability approaches discussed in Assignment 1.}
\end{table}

% ══════════════════════════════════════════════════════════════
\section{Conclusion}
% ══════════════════════════════════════════════════════════════

This report covered the theoretical foundations of propositional
satisfiability and the DPLL algorithm. The key findings are:

\begin{itemize}
  \item CNF is the universal input format for DPLL and modern SAT solvers;
        any propositional formula can be systematically converted to it.
  \item DPLL uses Unit Propagation, Pure Literal Elimination, and systematic
        backtracking to decide SAT/UNSAT completely and correctly.
  \item The formula $p \land \lnot\lnot(\lnot q \land \lnot\lnot p)$
        simplifies to $p \land \lnot q$ in CNF, satisfiable with
        $p = \top,\; q = \bot$.
  \item Z3 extends DPLL/CDCL to full SMT solving and confirmed the SAT result,
        returning the correct model $[p = \textit{True},\; q = \textit{False}]$.
\end{itemize}

\begin{tcolorbox}[greenstyle]
All the files are submitted as a zip file 
\end{tcolorbox}

\end{document}